% =============================================================================
% Chapter 4 -- Implementation
% =============================================================================
\chapter{Implementation}
\label{ch:impl}

This chapter documents the implementation choices that are
specific to this project rather than inherited unchanged from
the upstream COOM repository \citep{tomilin2024coom}. Three
choices were load-bearing: the Apple Silicon CPU TensorFlow
build, the sequential overnight experiment queue, and the
log-only reproducibility layer. Each is documented in turn. The
SAC backbone, the four regularisation methods and the
ViZDoom-Gymnasium glue are reused from the COOM repository
unmodified, beyond a single version pin on the Gymnasium
dependency that is recorded in the project history.

\section{Technical Stack}
\label{sec:stack}

The implementation builds on the COOM benchmark repository
\citep{tomilin2024coom} and uses TensorFlow~2.11
\citep{abadi2016tensorflow} with the
\texttt{tensorflow-probability} extension for the SAC actor.
ViZDoom~1.3 \citep{kempka2016vizdoom} provides the simulation
environment and is wrapped by COOM's own
\texttt{ContinualLearningEnv} class to expose tasks as a sequence
of Gymnasium-compatible environments.

The single non-trivial choice in this stack is the use of
TensorFlow's CPU build on Apple Silicon. The published COOM
training code targets CUDA, and porting to Apple's Metal backend
(via the \texttt{tensorflow-metal} plugin) was rejected after a
short pilot showed Metal's mixed-precision kernels did not match
the float32 numerics of CUDA closely enough to inherit the COOM
hyperparameters without re-tuning. The CPU build is slower by an
order of magnitude but has reproducible numerics, which matters
more than throughput for a comparative study. Empirically a
single CO4 run at $20{,}000$ steps per task completes in
$50$--$220$ minutes on an M-series CPU; the variance is mostly
explained by the per-method cost of computing parameter
importance at task boundaries (MAS being the slowest and
Fine-Tuning the fastest).

Python is pinned to 3.10 and the full set of pip-installable
dependencies is captured in the repository's \texttt{setup.py}.
The two non-Python dependencies are ViZDoom (a Doom binary
shipped via PyPI) and the system OpenGL stack on macOS, both of
which are stable across the duration of the project.

\section{Sequential Experiment Queue}
\label{sec:queue}

To make the most of a single CPU host without GPU parallelism,
experiments are scheduled by a small shell wrapper that records
start time, end time, return code and wall-clock duration to a
status file. The wrapper is single-threaded by design: running
two SAC trainers concurrently on the same machine destabilised
the timing of the ViZDoom screen buffer reader during pilot
runs, and the comparison-relevant metric --- success rate as a
function of training steps, not wall-clock time --- has no
incentive to parallelise.

\begin{code}[language=bash, caption={Sequential CPU experiment
    queue (excerpt from \texttt{CL/scripts/cpu\_queue.sh}).
    Each invocation appends a single record to the status file
    on completion.},label={lst:queue}]
run() {
    local tag="$1"; shift
    local t0=$(date +%s)
    "$@"
    local rc=$?
    local t1=$(date +%s)
    echo "[$tag] rc=$rc duration=$((t1 - t0))s" >> "$STATUS"
}

run ft   python CL/run_cl.py --sequence CO4 --seed 1 \
    --steps_per_env 20000 --no_test --group_id ft_co4_20k_seed1
run ewc  python CL/run_cl.py --sequence CO4 --seed 1 \
    --steps_per_env 20000 --no_test \
    --cl_method ewc --cl_reg_coef 250 \
    --group_id ewc_co4_20k_seed1
\end{code}

The full overnight queue covers six runs:
\{ Fine-Tuning, L2, MAS, EWC \} on seed~1, plus EWC on seeds 2
and~3 to support the seed band reported in \Cref{tab:avg_perf}.
The recorded status file (\texttt{queue\_status.txt}) is the
ground truth for run completion: each line records the
identifier, the return code, and the wall-clock duration. A
return code other than zero in this file is the signal that a
run requires re-execution, and the analysis pipeline refuses to
read logs from any run that is not present in the file with
return code zero.

\Cref{tab:walltime} reports the wall-clock cost of each method
at this sequence length and budget on the reference Apple
Silicon CPU host. The four-fold spread between the cheapest
method (L2 at $45$~min) and the most expensive (MAS at
$220$~min) is dominated by the per-task-boundary importance
computation, not by SAC training itself; this is consistent
with the design of MAS, which iterates a sample budget over
the model output gradient norm at end-of-task. The wall-clock
table is also the load-bearing evidence that the overnight
queue (\Cref{nfr:queue}) is feasible only at $20{,}000$ steps
per task: at the published $200{,}000$-step budget, MAS alone
would consume more than a day of CPU time per seed.

\begin{table}[t]
    \centering
    \caption{Wall-clock time per CO4 run at $20{,}000$ steps
        per task, on the reference Apple Silicon CPU host.
        Single-seed values report the run; multi-seed values
        report mean $\pm$ standard deviation.}
    \label{tab:walltime}
    \IfFileExists{tab4_walltime.tex}{%
        \begin{tabular}{lcc}
\toprule
Method & Wall-clock per CO4 run (min) & Seeds \\
\midrule
FT & 50 & 1 \\
L2 & 45 & 1 \\
MAS & 220 & 1 \\
EWC & 128 $\pm$ 8 & 3 \\
EWC-Online & 48 & 1 \\
\bottomrule
\end{tabular}%
    }{%
        \fbox{\parbox{0.6\linewidth}{\centering\vspace{1cm}%
            \textit{tab4\_walltime.tex placeholder.}%
            \vspace{1cm}}}%
    }
\end{table}

\section{Online EWC: An Improved Variant}
\label{sec:ewc_online_impl}

The diagnosis that the failure of EWC at the reduced budget is
driven by a Fisher matrix that captures source-task optimisation
noise (\Cref{sec:fishernoise}) suggests a single, minimally
invasive fix. The canonical formulation
\citep{kirkpatrick2017ewc} accumulates the per-task Fisher
diagonals additively across tasks,
\begin{equation}
    F^{(t)}_{\mathrm{total}} = \sum_{k=0}^{t} F^{(k)},
    \label{eq:fisher_sum}
\end{equation}
so the regularisation strength grows monotonically with the
number of tasks seen. Schwarz et al.\
(\citeyear{schwarz2018progress}) proposed instead to keep an
exponential moving average of the Fisher diagonal:
\begin{equation}
    F^{(t)}_{\mathrm{total}}
    = \gamma \cdot F^{(t-1)}_{\mathrm{total}}
      + (1 - \gamma) \cdot F^{(t)}, \quad \gamma \in [0, 1].
    \label{eq:fisher_ema}
\end{equation}
Two consequences follow. First, $F^{(t)}_{\mathrm{total}}$ is
bounded as the sequence grows, so the regularisation tax
ceases to scale with the number of tasks. Second, the EMA acts
as a low-pass filter on per-task Fisher noise: a single
unusually noisy Fisher estimate at the end of an
under-converged source task contributes only a fraction
$1 - \gamma$ of its magnitude to the running estimate, rather
than the full magnitude under \cref{eq:fisher_sum}. Both
consequences are exactly what the Fisher-on-noise diagnosis
predicts as the right intervention.

The implementation is one method override on the existing
EWC class. The new file \texttt{CL/methods/ewc\_online.py}
subclasses \texttt{EWC\_SAC} and replaces only the
Fisher-merge step:

\begin{code}[language=Python, caption={Online EWC: an
    exponential moving average replaces the additive Fisher
    accumulation in the original \texttt{EWC\_SAC}.
    Excerpt from \texttt{CL/methods/ewc\_online.py}.},
    label={lst:ewconline}]
class EWCOnline_SAC(EWC_SAC):
    def __init__(self, cl_reg_coef=1.0, regularize_critic=False,
                 ewc_gamma=0.95, **kwargs):
        super().__init__(cl_reg_coef=cl_reg_coef,
                         regularize_critic=regularize_critic,
                         **kwargs)
        self.ewc_gamma = float(ewc_gamma)
        self._fisher_initialised = False

    def _merge_weights(self, new_weights):
        if not self._fisher_initialised:
            # First merge: install F_1 directly via the same
            # broadcast addition the baseline uses.
            merged = list(old + new for old, new in
                          zip(self.reg_weights, new_weights))
            self._fisher_initialised = True
        else:
            g = self.ewc_gamma
            merged = list(g * old + (1.0 - g) * new
                          for old, new in
                          zip(self.reg_weights, new_weights))
        for old_w, new_w in zip(self.reg_weights, merged):
            old_w.assign(new_w)
\end{code}

The class is registered in the \texttt{CLMethod} enum of
\texttt{CL/run\_cl.py} and exposed through a single new
command-line flag \texttt{-{}-ewc\_gamma}, defaulting to $0.95$
on the recommendation of \citet{schwarz2018progress}. All
other hyperparameters are inherited from the original EWC
configuration verbatim, so the only difference between an
EWC run and an Online EWC run is the merge rule of
\cref{eq:fisher_sum} versus \cref{eq:fisher_ema}.

\section{Reproducibility Layer}
\label{sec:reprod}

Every figure and table in this dissertation is regenerated from
the recorded training logs by a single Python entry point. The
entry point reads only \texttt{progress.tsv} and
\texttt{config.json} files from each run directory, with no
dependency on TensorFlow checkpoints or in-process state. This
satisfies requirement~\ref{req:reproduce} and means a marker can
verify the headline numbers without re-running the training
queue.

The TSV loader handles a quirk of the upstream COOM code: as
training progresses, additional metric columns can be appended
to the right-hand side of \texttt{progress.tsv} but the header
line is written once at the start. The loader pads the column
list with synthetic names before the read, then projects back to
the original schema:

\begin{code}[language=Python, caption={Loader for COOM TSV
    training logs, handling the trailing-tab schema growth.
    See \Cref{lst:loader} in the appendix for the full
    implementation.},label={lst:loadersnippet}]
def load_run(group_dir):
    path = latest_progress(group_dir)
    with path.open() as f:
        header = f.readline().rstrip('\n').split('\t')
    names = header + [f"_pad{i}" for i in range(20)]
    df = pd.read_csv(path, sep='\t', skiprows=1, names=names,
                     engine='c', index_col=False)
    return df[header]
\end{code}

\section{Dependency Pinning}
\label{sec:pinning}

Because the CPU TensorFlow build is the load-bearing platform
choice (\Cref{sec:stack}), every Python dependency that
participates in a numerical computation is pinned. The pins live
in \texttt{setup.py} at the repository root; the values used to
produce the results in this dissertation were
TensorFlow~$2.11.0$, \texttt{tensorflow-probability}~$0.19.0$,
ViZDoom~$1.3.0$, \texttt{gymnasium}~$0.28.1$, and NumPy~$1.23.5$.
The \texttt{gymnasium} pin in particular is a project-specific
addition: the upstream COOM repository was originally written
against the now-deprecated \texttt{gym} package, and a small
shim was required to keep the existing environment registration
calls working under \texttt{gymnasium}. The pin was committed in
the \texttt{Specified gymnasium version} commit
on the project's \texttt{main} branch.

\section{Logging and Naming}
\label{sec:logging}

Each training run writes its outputs into a \texttt{group\_id}
directory named after the method, sequence, budget and seed, for
example \texttt{ewc\_co4\_20k\_seed1}. Within that directory a
timestamped sub-directory holds the actual logs, so re-running a
configuration appends rather than overwrites. This convention is
inherited from the Continual World provenance of the COOM
repository \citep{wolczyk2021cwsac} and is preserved here
without modification, because changing it would invalidate the
existing analysis scripts.

The Weights \& Biases \citep{wandb} integration shipped with the
upstream code is left in place but disabled by default in the
overnight queue: experiment tracking on a metered home internet
connection is not worth the bandwidth, and the local
\texttt{progress.tsv} files are sufficient for every figure
and table in this dissertation.
